\subsection{The full modular identity and a fourth-moment projection}

\begin{lemma}\label{lem:full-induction-sts}
Let $\mathcal C$ be a pseudo-unitary fusion category with commutative
Grothendieck ring, and let $b_\chi$ be the forgetful vector of the
actual summand $A_\chi\subset I(1)$ corresponding to the character
$\chi$. For a simple $Z\in Z(\mathcal C)$ write $v_Z$ for its
forgetful vector. Then
\begin{equation}\label{eq:full-induction-sts}
 \sum_Z\theta_Z(v_Z)_i(v_Z)_j
 =\sum_\chi(b_\chi)_i\overline{\chi(X_j)}.
\end{equation}
\end{lemma}
\begin{proof}
Let $C_{\chi i}=\chi(X_i)$, let $f_\chi$ be the formal codegrees,
and put $X=\operatorname{diag}(f_\chi^{-1})C$.
The character identification in \cite[Example~2.9 and
Theorem~2.13]{OstrikPiv3} gives
$S_{A_\chi,Z}=\chi(FZ)/f_\chi$ for every central simple $Z$,
with $S_{1Z}=\dim(Z)/\dim(\mathcal C)$.
All $A_\chi$ have twist one, and the center has trivial anomaly.
Thus the low block of $STS=T^{-1}ST^{-1}$ is
\[
 X\left(\sum_Z\theta_Zv_Zv_Z^{\mathsf t}\right)X^{\mathsf t}
 =XB^{\mathsf t},
\]
where $B$ has rows $b_\chi$. Character orthogonality gives
$C\overline C^{\mathsf t}=\operatorname{diag}(f_\chi)$,
so $X^{-\mathsf t}=\overline C$. Cancellation proves the identity.
\end{proof}

\begin{theorem}\label{thm:r18-sts-moments}
The ring $R8\_1\_18$ has no complex fusion categorification.
\end{theorem}
\begin{proof}
By \cite[Lemma~3.1]{SchNear}, pass to a pseudo-unitary Galois
conjugate. Put $d=2+\sqrt5$, $K=\mathbb Q(\sqrt5)$ and
$D=8+16d$. In model order
$(1,a,a^2,a^3,\rho,a\rho,a^2\rho,a^3\rho)$, mapped to source
$(0,1,7,2,3,6,5,4)$, the ring is commutative and satisfies
\[
 a^4=1,\qquad \rho^2=a+(1+a+a^2+a^3)\rho.
\]
Its six non-FP characters have formal codegree eight and are
\[
 \chi_{k,\pm}(a)=\zeta_8^{2k},\qquad
 \chi_{k,\pm}(\rho)=\pm\zeta_8^k,\qquad k=1,2,3.
\]
The other two codegrees are $D,D'$.
The two forced low rows are $1$ and $1+X_3+X_4+X_5+X_6$;
the other six have form $1+$ two nonpointed constituents.
Actual character reciprocity and duality leave eight low assignments.
Taking normalized relative trace $t_K$ in
\eqref{eq:full-induction-sts} gives, for nonunit indices,
\[
 \sum_Z(v_Z)_i(v_Z)_j t_K(\theta_Z)
 =\begin{cases}4-d&i,j\in\{3,4,5,6\},\\0&\text{otherwise}.
 \end{cases}
\]
Indeed every nonpointed value of the six $\mathbb Q(\zeta_8)$
characters has relative trace zero; the pointed entries cancel in
each of the eight assignments.

The pointed induction Gram is $8I_4$.
Lemma~\ref{lem:cyclic-induction-phase-mass} therefore gives
$(v_Z)_1,(v_Z)_2\leq1$ exactly as in the preceding theorem,
uniformly over all four pointed associator classes.
Besides the quadratic transport $T$ displayed there, full Galois
symmetry acts on rows by the source permutation
\[
 P=(0,2,1,6,5,4,3,7).
\]
It is induced by an automorphism fixing $K$ and sending
$\zeta_8$ to $\zeta_8^3$: on the model basis it sends
$a\mapsto a^{-1}$ and $a^j\rho\mapsto a^{3j+1}\rho$.
Include also source duality
$(0,2,1,4,3,6,5,7)$. Every full row orbit has a common
relative moment histogram: Galois transport changes a twist by
$\sigma^2$, which fixes quadratic $K$, and duality preserves twists.

The Cauchy theorem for spherical fusion categories
\cite[Theorem~3.9]{BNRWCauchy} applies to the nonintegral
dimension $D$. Its norm is $320=2^6\cdot5$, so the exponent has
only primes two and five. For powers $1,2,4$, every normalized
relative twist triple is one of the triples from fortieth roots
or the zero triple. This follows directly from the quadratic
Gauss sums: fifth-adic exponent at least two, or $2$-adic exponent
at least four, gives zero in all three coordinates. There are
exactly thirteen triples; this is not an assertion that the
exponent divides forty.

There is no need to determine fourth indicators. Each $\nu_4(X_i)$
is a sum of fourth roots, so $t_K(\nu_4(X_i))\in\mathbb Z$.
The functional $\ell(a+bd)=2a-b$ annihilates $D\mathbb Q$.
Applying it to the fourth induction identity gives
\[
 \ell\!\left(\sum_Z(v_Z)_i\dim(Z)t_K(\theta_Z^4)\right)=0.
\]
Together with the first two induction traces, the full modular
identity above, and the Gram identity, this gives $119$ rational
linear coordinates after low terms are subtracted.

The exact domain uses the diagonal Gram bounds, integrality of
$D/\dim(Z)$, the Fourier multiplicity bound, full row-orbit
closure, and positive semidefiniteness after subtracting that
entire orbit. It has $25$ row orbits for every low assignment.
The pointed classes $k=0$, $k=1,3$ and $k=2$ impose their respective
cyclic twist constraints and leave $167$, $145$ and $153$ columns.
The three integer covectors in
\texttt{r18\_sts\_moment\_certificate.json} pair nonnegatively
with every column and pair to $-256$ with every corresponding
target, for all eight low assignments. The standard-library
verifier reconstructs the source, characters, thirteen twist
triples and all bounded domains using exact rational arithmetic.
Thus none of the four pointed classes is possible. The initial
Galois reduction proves the general complex assertion.
\end{proof}
